\documentclass[a4paper,11pt]{article}
\usepackage[margin=27mm]{geometry}
\usepackage{fontspec,amsmath,amssymb,amsthm,unicode-math,booktabs}
\setmainfont{STIXGeneral}\setmathfont{STIX Two Math}
\newcommand{\gen}[1]{\langle#1\rangle}\newcommand{\Aut}{\operatorname{Aut}}
\newcommand{\mix}{\mathbin{\diamond}}\newcommand{\Alt}{\operatorname{Alt}}
\newtheorem{theorem}{Theorem}\newtheorem{proposition}[theorem]{Proposition}\newtheorem{corollary}[theorem]{Corollary}
\title{Polytopality of mixes and their regular covers}
\author{Human Compression project}\date{Author edition; pending independent review}
\begin{document}\maketitle
\begin{abstract}
We apply the common fiber-intersection argument to the criteria of
Cunningham--Hubard, in the complete author edition arXiv:2506.19334v1.
The main criterion uses every flag and is proved without a finiteness
assumption. Facet coverings, parity doubles and two-orbit regular covers
follow from the same test. We also exhibit a counterexample to the
printed coprime-label sufficient condition: the eleven-cell and the
four-simplex have a nonpolytopal mix. Several printed single-base conclusions require correction. A second
explicit construction explains this failure, while the all-flags theorem
and its precisely rooted consequences remain valid. This working edition
is awaiting independent scientific review.
\end{abstract}

\section{The main criterion and its premises}
Use the notation of the common paper. For two rooted components $R_1,R_2$
with flag stabilizers $N_1,N_2$ in their universal connection group, write
\[
 X(R_2\mid R_1)=N_2N_1/N_2
\]
when $N_1$ normalizes $N_2$. This is the \emph{lower} variance group, with
the order of the arguments shown explicitly. It is not merely an abstract
group type: for the facet, vertex-figure and medial groups below we use
the embeddings into $\Aut(P_2)$ proved in common Lemma~4.

\begin{theorem}\label{main}
Let $P_1$ be an $n$-polytope and $P_2$ an $n$-premaniplex. Suppose
both are admissible over a common rooted quotient $T$. For a flag
$q=(\Psi_1,\Psi_2)$ of their mix let $K_i,L_i,N_i$ denote its rooted
facet, vertex-figure and medial section in the $i$th factor.
Then $P_1\mix P_2$ is polytopal if and only if its facets and
vertex-figures are polytopal and, at every such flag,
\[
 X(K_2\mid K_1)\cap X(L_2\mid L_1)\leq X(N_2\mid N_1). \tag{1}
\]
\end{theorem}
\begin{proof}
The common quotient supplies a subgroup normalizing and containing both
flag stabilizers. Common Lemma~4 identifies the three variance groups
with the three fibers in common (6). Common Corollary~3 then gives the
assertion. All identifications are made at the specified flag. This is
[CH, Theorem~6.3], including the necessity in its Proposition~6.2.
\end{proof}

The restricted facets and vertex-figures are themselves rooted mixes.
Thus when all relevant rooted facet mixes and vertex-figure mixes are
polytopal, a trivial intersection on the left of (1) suffices, giving
the form of [CH, Corollary~6.4] with its flag choices explicit.
Common Proposition~6 proves [CH, Proposition~6.5]: a facet covering
at every paired flag forces polytopality. Its proof includes the required
induction on the vertex-figures, not just vanishing of one variance group.

When the common quotient remains connected on middle colors, common
Proposition~5 permits one base flag in (1). More generally one needs only
one representative per middle-color component of that quotient. In
particular there is one test for regular inputs and for two-orbit inputs
whose common two-vertex quotient has an edge of an interior color.
These are conclusions of an actual common symmetry action on the mix.
We have not replaced that hypothesis by separate transitivity assertions
about its two factors.

\section{Why separate medial transitivity does not give one test}
The common-quotient condition in the preceding section cannot simply be
replaced by medial-section transitivity of the two inputs separately.
We give two isomorphic polytopes with different compatible rootings that
satisfy all the local and symmetry hypotheses in [CH, Corollaries~6.4,
6.6, Theorem~6.7 and Proposition~6.8]. The indicated base test passes,
although their mix is not polytopal. In particular, this is not a
counterexample depending on nonuniform facets or on a nonpolytopal input.

\subsection{A six-coordinate action}
Let $S=S_5$, with $s_i=(i+1\ i+2)$ for $0\leq i\leq3$.
Its six Sylow $5$-subgroups are indexed $0,\ldots,5$ using the respective
cycle generators
\[
 (12345),\ (12354),\ (12453),\ (12435),\ (12543),\ (12534).
\]
Conjugation induces a permutation action $\pi$ on these six coordinates.
Writing each image list in order $0,\ldots,5$, the four permutations are
\[
\begin{array}{c|rrrrrr}
i&\pi(s_i)(0)&\pi(s_i)(1)&\pi(s_i)(2)&\pi(s_i)(3)&\pi(s_i)(4)&\pi(s_i)(5)\\\hline
0&4&2&1&5&0&3\\
1&5&3&4&1&2&0\\
2&3&2&1&0&5&4\\
3&1&0&4&5&2&3
\end{array}                                                    \tag{S1}
\]
Use the induced left permutation action on
$V=\{v\in\mathbb F_3^6:\sum_hv_h=0\}$, so $(gv)_h=v_{\pi(g)^{-1}h}$.
In $V\rtimes S$ put $R_i=(u_i,s_i)$, where
\[
 u_0=u_2=u_3=0,\qquad u_1=(-1,-1,0,1,0,1).
\]
Multiplication is $(u,g)(v,h)=(u+gv,gh)$. Direct substitution in (S1)
shows $R_i^2=1$ and $[R_i,R_j]=1$ for $|i-j|>1$; also
\[
 (R_0R_1)^3=(R_2R_3)^3=1,\qquad
 c:=(R_1R_2)^3=(v,1),\quad v=(1,1,-1,-1,1,-1).       \tag{S2}
\]
These identities are small coordinate calculations, all displayed inputs
being fixed elements of $\mathbb F_3$ and $S_5$.

The four marks generate all of $G=V\rtimes S$. Indeed their projection
is $S$, and their normal closure of $c$ contains the $S$-span of $v$.
The five vectors $v,s_0v,s_3v,s_0s_3v,s_1s_0v$ have rank five over
$\mathbb F_3$; their rows are
\[
\begin{pmatrix}
1&1&2&2&1&2\\1&2&1&2&1&2\\1&1&1&2&2&2\\
2&1&1&2&1&2\\2&2&1&2&1&1
\end{pmatrix}.
\]
They lie in $V$, so span $V$. Consequently $|G|=3^5\cdot120=29160$.

For $h\in\{0,\ldots,5\}$ let $N_h=\{v\in V:v_h=0\}$ and form the
colored Schreier graph $P_h=N_h\backslash G$, with adjacency given by
right multiplication by $R_i$. Using right rather than left connection
actions merely reverses the names of words; the colored path statements
are unchanged. A flag can be written $(g,k)\in S\times\mathbb F_3$, and
its $i$-adjacency is
\[
 (g,k)^i=(gs_i,k+(g u_i)_h).                            \tag{S3}
\]
Thus every $P_h$ has $360$ flags and covers the simplex flag graph
$T=S$. The cover is regular with deck group $V/N_h\cong C_3$.
Equivalently, the preimages of $N_h\trianglelefteq V$ in the universal
connection group satisfy exactly the source's definition of
$T$-admissibility. Since $T$ is a maniplex, so are these covers.

\subsection{Every input is a medial-section-transitive polytope}
Write $A=\gen{R_0,R_1,R_2}$, $B=\gen{R_1,R_2,R_3}$ and
$C=\gen{R_1,R_2}$. Their translation kernels under $G\to S$ are
\[
\begin{aligned}
 V_A&=\gen{e_0-e_5,e_1-e_2,e_3-e_4}_{\mathbb F_3},\\
 V_B&=\gen{e_0-e_3,e_1-e_5,e_2-e_4}_{\mathbb F_3},\\
 V_C&=\gen{v}_{\mathbb F_3}.
\end{aligned} \tag{S4}
\]
Here is a full justification of these kernel formulas. The projected
end groups are $S_4$, with their usual adjacent-transposition
presentations. All its Coxeter relations already hold in each lifted
end except the middle cube in (S2). Therefore its kernel is generated
by the conjugates of that cube, and equals the linear span of the
corresponding orbit of $v$. The first space in (S4) is invariant under
$s_0,s_1,s_2$, contains $v$, and the three vectors $v,s_0v,s_1s_0v$
are independent in it. The second space is invariant under
$s_1,s_2,s_3$, and contains the independent vectors $v,s_3v,s_2s_3v$.
This proves both equalities. The last follows also from the presentation
of a group generated by two involutions: $R_1R_2$ has order nine and
its projection has order three. We retain here only the elementary
Coxeter presentations of symmetric groups; the translation kernels are
computed in full above.

The restricted groups on colors $\{0,1\}$ and $\{2,3\}$ project
isomorphically to $S_3$, by (S2). Consequently each facet of $P_h$ covers
a tetrahedron with trivial fiber on its $\{0,1\}$ sections. The common
fiber criterion gives the rank-three intersection property: one of the
two fibers being intersected is trivial. Its two-color sections are
polygons, so every facet is polytopal. The dual argument proves the same
for every vertex-figure. Each facet and vertex-figure has $24\cdot3=72$
flags, since every coordinate projection of both spaces in (S4) is
surjective. Every medial section has $6\cdot3=18$ flags, since all
coordinates of $v$ are nonzero. At every flag of $P_h$ the medial fiber
over $T$ is therefore the whole $C_3$ fiber. The common criterion now
proves that $P_h$ itself is polytopal. This proves all-flag assertions,
not only a test at a selected root.

The stabilizer $H_h$ of the $h$th Sylow subgroup in $S_5$ is its
normalizer, an affine group of order $20$ on five points. To see this
directly, identify the five points with $\mathbb F_5$: translations
normalize the chosen regular $C_5$, as do the four nonzero scalings,
and a normalizer is determined by the image of one point and of a
cycle generator. This group is sharply two-transitive. Since
$V\rtimes H_h$ normalizes $N_h$, its left action on $N_h\backslash G$
commutes with all colored adjacencies. On $T$, a middle-color component
is specified by the ordered pair $(g(1),g(5))$. There are twenty such
pairs, and $H_h$ is transitive on them. Each of these components has
one connected preimage in $P_h$, as just proved. Hence $P_h$ is
medial-section-transitive. In particular its facets, vertex-figures and
medial sections have uniform isomorphism types. The medial components themselves form one
automorphism orbit, exactly as required by the source. The six inputs
are isomorphic: conjugation by a suitable $(0,a)$ sends $N_h$ to $N_k$,
and left multiplication by $(0,a)$ then gives a color-preserving map
$N_h\backslash G\to N_k\backslash G$.

\subsection{One base passes and another fails}
Consider $P_0$ and $P_1$, rooted at compatible flags specified below.
Their mix is $(N_0\cap N_1)\backslash G$, with $1080$ flags. Its flags
are $(g,k,l)$, and (S3) holds in both sheet coordinates. All its facets
and vertex-figures are polytopal, since they are mixes of polyhedra.

At a flag above $g\in S$, the three fibers of this mix over $T$ are the
images of $gV_A,gV_B,gV_C$ under the two coordinate projections
$(v_0,v_1)$. This follows by traversing a word whose projection returns
to $g$ in the corresponding color subgroup. At $g=1$, (S4) gives
\[
 D_A=D_B=\mathbb F_3^2,\qquad D_C=\gen{(1,1)}.
\]
Thus the end intersection fails: it has $54$ flags whereas the middle
component has $18$. There is also an explicit path witness. Put
$h=R_2R_3$ and use the translation $c$ from (S2). The words
\[
 (cR_0cR_0)^2\in A,\qquad (chc^{-1}h^{-1})^2\in B
\]
both carry $(1,0,0)$ to $(1,1,0)$, by (S1)--(S2). No middle-color word
can do so, since the middle translation fiber is the diagonal line.
The mix is consequently not polytopal.

Now set $g_*=s_1s_2s_3=(2\,3\,4\,5)$, and use base flags with
$(g,k,l)=(g_*,0,0)$. From (S1), $\pi(g_*)=(4,1,0,2,3,5)$.
The coordinate pair $(0,1)$ of $g_*z$ is therefore $(z_2,z_1)$.
Equations (S4) give at this base
\[
 D_A=\gen{(-1,1)},\qquad D_B=\mathbb F_3^2,
 \qquad D_C=\gen{(-1,1)}.                              \tag{S5}
\]
The base intersection thus holds. To express the test exactly as a mix
of $P_0$ with $P_1$, rather than a cover of $T$, require that the first
sheet remain fixed. Intersect the three fibers in (S5) with
$\{0\}\times\mathbb F_3$. The facet and middle fibers become trivial;
the vertex fiber is $C_3$. In particular the two end variance groups
have trivial intersection at the specified base.

Here and throughout this paper, medial-section transitivity means
transitivity on the middle-color components, as used in the proof of
[CH, Proposition~2.5] and its two-orbit tests. The parenthetical chain
type on the source's physical page~6 should instead be $\{0,n-1\}$;
it prints the complementary type. We do not silently conflate these
notions.

Both inputs are polytopes and individually medial-section-transitive,
both are admissible over the same regular $T$, all their local types
are uniform, and even every mixed end is polytopal. Nevertheless the
base conditions hold while the mix fails elsewhere. This disproves the
single-base conclusions of [CH, Corollary~6.4 and Theorem~6.7] as
printed, and of Corollary~6.6 and Proposition~6.8 under the operational
medial-section convention just specified. The source's all-flags Theorem~6.3
is unaffected. A correct replacement either keeps every flag, assumes
medial-section transitivity of the mix itself, or uses the explicitly
proved common-quotient condition $W=CV$. The same argument identifies
exactly what went wrong: separate automorphisms of the two factors need
not act together on their chosen mix.


\section{Polyhedra and coprime middle labels}
The mix of any two rooted polygons is a polygon: its connected flag graph
is again a simple two-colored cycle or double ray. In fact the mix of any
two polyhedra is polytopal, without a symmetry assumption. To prove this,
apply the common fiber criterion at colors $\{0,1\}$ and $\{1,2\}$.
If a point of the common fiber belongs to both, projection to the second
polyhedron places it either at $y$ or at $r_1y$. The latter is impossible.
A word in $r_0,r_1$ fixing the first polygon flag has even length, whereas
a word carrying the second polygon flag to its $1$-adjacent flag has odd
length. Both assertions follow from the bipartite flag graph of a
polygon. Thus the fiber intersection is trivial, as required. This
supplies the input cited from earlier work in [CH, Proposition~6.8].
For rank four, therefore, condition (1) is the only remaining test once
the appropriate flag representatives have been justified.

\begin{proposition}\label{coprime}
Let $P_i$ be orientable regular four-polytopes of types
$\{p_i,q_i,r_i\}$, where $q_1,q_2$ are finite and coprime.
Then $P_1\mix P_2$ is polytopal.
\end{proposition}
\begin{proof}
In the second medial dihedral group $D_{2q_2}$, the lower variance group
of a $q_2$-gon relative to a $q_1$-gon is the full rotation subgroup:
$\gen{(\rho_1\rho_2)^{q_1}}=\gen{\rho_1\rho_2}$.
The intersection of the two end variance groups lies in
$\gen{\rho_0,\rho_1,\rho_2}\cap
\gen{\rho_1,\rho_2,\rho_3}=\gen{\rho_1,\rho_2}$,
because $P_2$ is polytopal. Every word fixing a first-factor facet flag
has even length, since $P_1$ is orientable. Its image in the second group
is therefore even under the orientation character of $P_2$.
Consequently the common variance group lies in the rotation subgroup of
the medial dihedral group. Theorem~\ref{main} and the polyhedron result
prove the claim. It is enough in this proof that one first-factor end
be orientable, while the second polytope is orientable.
\end{proof}

The orientability assumption on the first factor cannot be dropped as
it is in the literal [CH, Proposition~6.9]. The following is a complete
finite obstruction, using the familiar eleven-cell. Its marked group
and intersection property are already established in our projective-group
paper; a short standalone finite certificate is being bound here so that
the cost of that input is visible.

Let $P_1$ have the following involutions on eleven points:
\[
\begin{aligned}
 a_0&=(1\,2)(3\,5)(4\,7)(6\,8),\\
 a_1&=(1\,3)(2\,4)(5\,7)(6\,9),\\
 a_2&=(3\,6)(4\,7)(5\,8)(9\,10),\\
 a_3&=(1\,3)(2\,5)(4\,7)(10\,11).
\end{aligned}                                                    \tag{2}
\]
These generate the simple group of order $660$, with a C-marking of type
$\{3,5,3\}$. Both end groups are $A_5$ and their intersection is the
middle dihedral group of order ten. Let $P_2$ be the four-simplex, with
$b_i=(i+1\ i+2)$ in $S_5$. It is orientable, via the sign character, and
its middle label is three, coprime to five.

The mixed facet projects onto $A_5$ and $S_4$. A subdirect product of
these groups is their full direct product: by the elementary common
quotient description of a subdirect product, a proper one would give a
nontrivial common quotient; simplicity of $A_5$ and $60>24$ exclude this.
The same argument applies to the other end. Their intersection in the
product group is consequently
\[
 D_{10}\times S_3,
\]
of order $60$. The common marked middle is generated by two involutions
whose product has order $\operatorname{lcm}(5,3)=15$, so is dihedral of
order $30$. More explicitly, $(a_1,1)$ belongs to both ends. It does not
belong to the common marked middle, since in that middle the reflection
character of $D_{10}$ equals the sign character of $S_3$, whereas their
values on $(a_1,1)$ differ. Thus the mix fails the intersection property.
Both inputs are regular polytopes, the second is orientable, and the
middle labels are coprime, exactly as required in the printed proposition.
The defect in its proof is that evenness of the first-factor relations
was inferred from orientability of the second factor.

\section{A common quotient faithful on a medial section}
Suppose $P_1,P_2$ are regular $n$-polytopes with facets $K_1,K_2$.
Suppose both facets cover a regular premaniplex $K_3$, and the covering
from $K_2$ is an isomorphism on its vertex-figures. This is equivalent
to the source's isomorphism-type hypothesis: two isomorphic regular
sections have a base-preserving isomorphism, obtained by composing with
an automorphism, and a covering preserving that base is uniquely
determined by colored paths. Thus the given covering is that isomorphism,
also for infinite sections. Then their mix is
polytopal, as in [CH, Proposition~6.14]. Here is the full kernel argument.
Every relation of $K_1$ maps trivially to $K_3$, so
$X(K_2\mid K_1)$ lies in the kernel of $\Aut(K_2)\to\Aut(K_3)$.
That kernel meets the base vertex-figure subgroup trivially. The
intersection with $X(L_2\mid L_1)$ lies in that very subgroup, by the
intersection property of $P_2$, and is consequently trivial.

There are still the local polytopality obligations. The vertex-figures
of $K_1$ cover those of $K_3$, hence those of $K_2$. The dual of the
common facet-cover proposition makes $K_1\mix K_2$ polytopal.
Likewise the facets of the vertex-figure $L_1$ cover those of $L_2$,
so $L_1\mix L_2$ is polytopal. Theorem~\ref{main} finishes the argument.
This proves the uniform statement; applying it to square or cubic toroids
retains the classical fact that their translation quotient with unit
period has the same spherical vertex-figure. We do not count the
classification of regular toroids as a new consequence of this proof.

\section{Two tests with their actual embeddings}
For completeness we give the exceptional two-orbit reduction without
suppressing the root of any variance group. This is the precise
interpretation needed when applying [CH, Corollaries~6.10--6.12].
Let $V=\ker\epsilon_I$ contain and normalize both input stabilizers,
and suppose every interior color belongs to $I$. Choose a paired base
flag $q_0$ and the second representative $q_1=r_jq_0$, where
$j\notin I$. Write
\[
 A_k=X(K_2(q_k)\mid K_1(q_k)),\quad
 B_k=X(L_2(q_k)\mid L_1(q_k)),\quad
 C_k=X(N_2(q_k)\mid N_1(q_k)).
\]
Each is embedded in the actual automorphism group of the second input
using its displayed paired root. The exact criterion is local
polytopality at the listed types and
\[
 A_k\cap B_k\leq C_k\qquad(k=0,1).                  \tag{9}
\]
The common orbit theorem proves sufficiency and necessity. The number
of local types and the possible common subgroup are as follows.
\[
\begin{array}{c|c|c|c|c}
 I^c&j&\text{facet types}&\text{vertex types}&\text{common subgroup}\\\hline
 \{0,n-1\}&0&1&1&A_0=A_1\\
 \{n-1\}&n-1&2&1&B_0=B_1\\
 \{0\}&0&1&2&A_0=A_1
\end{array}
\]
Indeed a facet orbit corresponds to a component of the two-vertex
quotient after deleting color $n-1$, and similarly for vertex orbits.
This proves the type counts, including coincident isomorphic types.
To justify the subgroup equalities, the global lower variance group
acts on the mix by automorphisms which fix its first coordinate. Its
subgroup stabilizing a mixed facet is exactly the facet variance group:
an element preserves that facet if and only if it sends its base flag
inside the same facet. The stabilizer is independent of the chosen
flag inside that facet. Thus moving the base by $r_0$ leaves $A_k$
unchanged. Moving it by $r_{n-1}$ similarly leaves $B_k$ unchanged.
This proves every equality in the table.

The last two rows reproduce the two displayed tests of Corollaries~6.11
and 6.12 with their embeddings explicit. In the first row, (9) reads
$A_0\cap B_0\leq C_0$ and $A_0\cap B_1\leq C_1$.
Even when both vertex-figures have the same isomorphism type, their
embeddings must be transported before identifying $B_0$ and $B_1$.
If an actual common embedding makes them equal, these two tests become
$A_0\cap B_0\leq C_0\cap C_1$, the expression printed in
Corollary~6.10. The table (9) is the unconditional replacement. We do
not infer equality of embedded subgroups solely from equality of
abstract rooted-section types, and do not claim a counterexample to
the literal two-orbit version of Corollary~6.10 here.


\section{Parity doubles}
Let $P=W/N$ be a regular polytope. Fix a nonzero character
$\epsilon:W\to C_2$ whose values on the generators prescribe the chosen
flag coloring, and let $Q$ be the mix of $P$ with the two-flag quotient
$W/\ker\epsilon$. Its stabilizer is $N\cap\ker\epsilon$.
For every color subgroup $S$ its lower variance subgroup in $C_2$ is
exactly
\[
 E_S=\epsilon(N\cap S).                                    \tag{3}
\]
It is trivial precisely when the coloring descends to that rooted section
of $P$. The character formulation also handles a restriction that is
zero on all its colors: the corresponding connected quotient section is
then a single flag.

Write $A,B,C$ for the facet, vertex-figure and middle color subgroups.
If $E_A=1$ or $E_B=1$, the corresponding facet or vertex-figure of $P$
covers that section of the two-flag quotient. The common facet-cover
proposition or its dual gives polytopality of $Q$.
If $E_A=E_B=C_2$ and $E_C=1$, (1) fails and $Q$ is not polytopal.
Finally, if $E_C=C_2$, (1) is automatic and $Q$ is polytopal precisely
when both mixed ends are polytopal. These are all three clauses of
[CH, Proposition~6.13], with restricted trivial characters made explicit.

\section{The smallest regular cover of a two-orbit polytope}
Let $P=W/N$ be a two-orbit polytope of class $2_I$. Put
$V=\ker\epsilon_I$, so $N$ is normal in $V$ and $[W:V]=2$.
For $s\in W\setminus V$, the normal core of $N$ in $W$ is
\[
 N\cap sNs^{-1}.                                          \tag{4}
\]
Indeed all conjugates under $V$ agree, and the remaining conjugates are
$sNs^{-1}$. Hence the smallest regular premaniplex cover is the mix of
$P$ with its other rooting $\overline P$. This is a statement about the
smallest regular \emph{premaniplex} cover; whether it is polytopal remains
the test (1), not an automatic consequence of (4).

Because (4) is normal in $W$, the mix $P\mix\overline P$ is regular.
The final observation in the common orbit theorem therefore gives a
single-base criterion for \emph{every} two-orbit class: its two rooted
mixed ends must be polytopal, and their two lower variance groups must
have intersection in the medial variance group at that same base.
The extra automorphisms of this regular mix may interchange the two
coordinate roles; the proof does not require them to project to
automorphisms of either input.

This supplies all four class-dependent decision criteria in the subject
of [CH, Corollary~6.15]. To identify the actual ends, choose the paired
roots $(\Phi,s\Phi)$ in the underlying polytope $P$. If $P$ is
facet-transitive, its two facets have the same unrooted type $K$, and
their rooted mix is $K\mix\overline K$. If its facets have two
orbits, the two roots pick the two types $K,K'$, giving $K\mix K'$.
The analogous alternatives give $L\mix\overline L$ or $L\mix L'$
for the vertex-figures. A facet has two orbits only for
$I^c=\{n-1\}$, and a vertex-figure has two orbits only for
$I^c=\{0\}$, by the quotient components in the preceding table.
The medial term is formed at the same paired root, so it is the
variance of the two specified medial sections, including distinct
types when $I^c\subseteq\{0,n-1\}$.

If one wishes to retain the source's two representative format, apply
(9) to these two re-rootings. Regularity of the mix makes the two
path tests equivalent, although the displayed coordinate embeddings
can differ. The source's simplified intersection involving two medial
variance groups must be read with those transport data. Our criterion
keeps one actual paired base and makes no additional identification.

Here is the further simplification behind [CH, Proposition~6.16 and
Corollary~6.17]. If $n-1\notin I$, each $(n-2)$-face is regular.
Given two flags of that face, extend them to flags of $P$. If they are
in different global automorphism orbits, change the second extension
in color $n-1$. This preserves the given lower face flag and changes
its global orbit. An automorphism now maps the two extensions and
restricts to the required face automorphism. If also $|I|<n-1$, the
facets of $P$ belong to one orbit: deleting the last color from the
two-vertex quotient leaves an edge. They therefore have one type $K$,
whose facets are regular. The mix $K\mix\overline K$ has identical
rooted regular facets in the two factors, so the common facet-cover
proposition proves that it is polytopal.

When both extreme colors lie outside $I$, this applies to both ends.
If in addition an interior color lies outside $I$, the one-base criterion
therefore reduces to
\[
 X(K\mid\overline K)\cap X(L\mid\overline L)
 \leq X(N_0\mid\overline{N_0}),
\]
where $N_0$ now denotes the medial section, not a flag stabilizer.
This recovers the chiral specialization and the stated two-orbit extension
with the medial-orbit hypothesis explicitly discharged. The borderline
case with no such interior color belongs to the two-test list above.

\section{Replacing the second factor by an admissible cover}
Suppose $N_1\trianglelefteq V$ and the second stabilizer $N_2$ is normal
in all of $W$. The mix $T\mix P_2$ has stabilizer $V\cap N_2$, normal
in $V$. Moreover
\[
 N_1\cap(V\cap N_2)=N_1\cap N_2.
\]
Thus $P_1\mix P_2=P_1\mix(T\mix P_2)$ as rooted premaniplexes,
and the latter is a mix of two $T$-admissible inputs. This proves the
replacement in [CH, Section~6.6]. If the first input is polytopal, the
main theorem applies without a polytopality requirement on the refined
second factor. If only the second is known polytopal, one must separately
establish polytopality of its refinement; the parity-double analysis
shows why that obligation is real.

\section{Current source boundary}
The common paper and the preceding sections give complete written proofs of the
all-flags criterion, the normalizer interpretation and the displayed
consequences. The single-base formulations are replaced by the proved orbit condition
and the explicit six-coordinate obstruction. The exceptional two-orbit
rows retain their actual embedded subgroups as in (9); the stronger
identification suppressed in Corollary~6.10 is not imported as a premise. The printed Proposition~6.9 has the explicit
obstruction above and a corrected orientable sufficient condition.

The three questions in [CH, Section~6.6] are not claimed solved. No
numerical factor or principal-source credit is attached to this working
edition. The complete common theory is part of any eventual proof and
measurement boundary, independently of its overlap with earlier topics.

\begin{thebibliography}{9}
\bibitem{CH} G. Cunningham and I. Hubard, \emph{Polytopality criteria for the
mix of polytopes and maniplexes}, arXiv:2506.19334v1, 24 June 2025,
21 physical pages.
\end{thebibliography}
\end{document}
